\documentclass[11pt]{article}
\usepackage[a4paper,margin=1in]{geometry}
\usepackage{amsmath,amssymb,amsthm}
\usepackage{booktabs}
\usepackage{array}
\usepackage{hyperref}

\newtheorem{definition}{Definition}
\newtheorem{theorem}{Theorem}

\newcommand{\Fq}{\mathbb F_q}
\newcommand{\Qell}{\mathbb Q_\ell}
\newcommand{\Xbar}{\overline X}
\newcommand{\XiWD}{\Xi_{\mathrm{WD}}}
\newcommand{\CRE}{\mathrm{CRE}}
\newcommand{\CRCFT}{\mathrm{CRCFT}}

\title{Step 186: Weil--Deligne Function-Field RH Carrier Pivot}
\author{RH Membrane Adequacy Track}
\date{}

\begin{document}
\maketitle

\section*{Verdict}

\[
\boxed{\mathrm{weil\_deligne\_verdict}=V_{\mathrm{weil\_deligne\_closes}}.}
\]

The Weil/Deligne function-field carrier is RH-analogous but not
classically-RH-equivalent for Riemann's zeta.  On its native cohomological
carrier, the framework residual closes:
\[
\boxed{\XiWD=0}
\]
by Grothendieck's trace formalism and Deligne's purity theorem.

\section*{Inherited Record Audit}

The inherited local records do not contain a full function-field
Weil/Deligne carrier package.  The only nearby inherited item is the
Deninger/motivic context recorded in the Connes pivot, which is explicitly
contextual and not a constructed number-field cohomology.  Thus this step
declares the function-field carrier directly and uses classical
Weil--Grothendieck--Deligne theorems as the native ledger.

\section*{Carrier Declaration}

Let \(X/\Fq\) be a smooth projective variety of dimension \(d\), and let
\(\ell\ne \operatorname{char}(\Fq)\).  Define the graded cohomological carrier
\[
H_{\mathrm{WD}}(X)
=
\bigoplus_{i=0}^{2d}
H^i_{\mathrm{\acute et}}(\Xbar,\Qell).
\]
This is a finite-dimensional graded \(\Qell\)-cohomological carrier with
Poincare duality pairings.  It is not a Hilbert carrier for Riemann's zeta.

The native operator is Frobenius:
\[
\mathrm{Frob}_q:
H^i_{\mathrm{\acute et}}(\Xbar,\Qell)
\longrightarrow
H^i_{\mathrm{\acute et}}(\Xbar,\Qell).
\]
We fix a geometric Frobenius convention; switching to arithmetic Frobenius
inverts the eigenvalues and must be tracked consistently.

\section*{Native Probes and Trace Ledger}

Native probes are cohomology classes and Frobenius-equivariant constructions.
Dissolving probes are the trace measurements
\[
\operatorname{Tr}\!\left(\mathrm{Frob}_q^n\mid
H^i_{\mathrm{\acute et}}(\Xbar,\Qell)\right).
\]
Grothendieck--Lefschetz gives the point-count ledger:
\[
\#X(\Fq^n)
=
\sum_{i=0}^{2d}(-1)^i
\operatorname{Tr}\!\left(\mathrm{Frob}_q^n\mid
H^i_{\mathrm{\acute et}}(\Xbar,\Qell)\right).
\]
The zeta function is
\[
Z(X,t)
=
\exp\left(\sum_{n\ge1}\#X(\Fq^n)\frac{t^n}{n}\right)
=
\prod_i
\det\!\left(1-t\,\mathrm{Frob}_q\mid
H^i_{\mathrm{\acute et}}(\Xbar,\Qell)\right)^{(-1)^{i+1}}.
\]
With \(t=q^{-s}\), an eigenvalue \(\alpha_{i,j}\) of Frobenius on \(H^i\)
contributes a zero or pole at \(\alpha_{i,j}q^{-s}=1\), with the sign
determined by the cohomological parity.

\section*{Parent Residual}

\begin{definition}[Weil--Deligne residual]
\(\XiWD\) is the graded Frobenius-weight defect:
\[
\XiWD=0
\quad\Longleftrightarrow\quad
|\iota(\alpha_{i,j})|=q^{i/2}
\]
for every Frobenius eigenvalue \(\alpha_{i,j}\) on
\(H^i_{\mathrm{\acute et}}(\Xbar,\Qell)\) and every complex embedding
\(\iota:\overline{\Q}_\ell\to\mathbb C\) used to measure absolute values.
\end{definition}

Equivalently, \(\XiWD=0\) says the native degree-\(i\) zero/pole ledger lies
on the expected line \(\operatorname{Re}(s)=i/2\), after the parity and
reciprocal conventions of \(Z(X,t)\) are accounted for.

\begin{theorem}[Native Weil--Deligne closure]
For smooth projective \(X/\Fq\),
\[
\XiWD=0.
\]
\end{theorem}

\begin{proof}
Grothendieck's \(\ell\)-adic cohomological formalism supplies the legal native
carrier, the Lefschetz trace formula, Poincare duality, and the determinant
factorization of \(Z(X,t)\).  Deligne's proof of the Weil conjectures gives
purity: Frobenius on \(H^i_{\mathrm{\acute et}}(\Xbar,\Qell)\) is pure of
weight \(i\).  Therefore every eigenvalue satisfies
\[
|\iota(\alpha_{i,j})|=q^{i/2}.
\]
This is exactly the vanishing condition for the graded weight defect
\(\XiWD\).  Thus the framework residual closes natively on the
function-field carrier.
\end{proof}

\section*{CRE and CRCFT Audit}

\begin{center}
\begin{tabular}{p{0.28\linewidth}p{0.18\linewidth}p{0.44\linewidth}}
\toprule
Audit & Result & Reason\\
\midrule
CRE status & not \(\CRE\)
& Closing \(\XiWD\) proves function-field RH for \(X/\Fq\), not Riemann RH.\\
\(\CRCFT\) applicability & does not apply
& \(\CRCFT\) is a foreclosure taxonomy for CRE carriers; this non-CRE carrier closes natively.\\
Framework closure & closed native
& Grothendieck--Lefschetz plus Deligne purity supplies the full cohomological ledger.\\
Transfer to Riemann RH & not claimed
& No number-field cohomology carrier with the same properties is supplied.\\
\bottomrule
\end{tabular}
\end{center}

\section*{Pattern Implication}

Together with the Selberg/Maass carrier, this gives a second non-CRE
RH-analogous carrier that closes natively.  The emerging distinction is now
sharper:
\[
\text{CRE carriers tested} \Rightarrow \CRCFT\text{ foreclosure modes,}
\]
while
\[
\text{proved non-CRE RH analogs} \Rightarrow \text{native closure when their full ledger is present.}
\]
This is an empirical framework pattern, not a proof that all future non-CRE
analogs close.

\section*{Classical Theorems Cited}

\begin{itemize}
\item A. Weil, 1948, \emph{Sur les courbes algébriques et les variétés qui
s'en déduisent}: function-field RH for curves and the geometric method.
\item A. Weil, 1949, \emph{Numbers of solutions of equations in finite
fields}: formulation of the Weil conjectures.
\item A. Grothendieck, 1965 and SGA: \(\ell\)-adic cohomology,
Grothendieck--Lefschetz trace formula, determinant factorization, and
functional equation formalism.
\item P. Deligne, 1973, \emph{La conjecture de Weil I}: proof of RH/purity
for smooth projective varieties over finite fields.
\item P. Deligne, 1980, \emph{La conjecture de Weil II}: weights and mixed
sheaf formalism.
\end{itemize}

\section*{Nonclaim Boundary}

This step does not prove Riemann RH, does not transfer function-field RH to
the number-field zeta function, and does not construct a number-field
cohomology with Deligne purity.  It closes only the native Weil/Deligne
function-field residual \(\XiWD\).

\end{document}
